\chapter{CONCLUSION AND FUTURE RECOMMENDATIONS}

\section{Conclusion}

This project presented the design and implementation of an AI-powered Academic Support System based on the Retrieval-Augmented Generation (RAG) architecture. The system was developed to assist students by generating syllabus-oriented academic notes from institution-specific learning materials. By combining semantic retrieval using the Qdrant vector database with the Google Gemini Large Language Model, the system produces contextually relevant and structured educational content while reducing reliance on the language model's internal knowledge. \par

The implemented system successfully performs document ingestion, semantic retrieval, and automated note generation through a modular architecture consisting of routers, services, adapters, and utility modules. The use of vector embeddings enables efficient retrieval of relevant academic content, while the integration of a Large Language Model provides clear and coherent explanations. Unit and system testing demonstrated that the implemented modules function correctly and meet the intended project objectives. \par

Overall, the proposed system provides a practical foundation for AI-assisted academic support and demonstrates the applicability of Retrieval-Augmented Generation techniques in educational environments. \par

\subsection{Future Recommendations}

Although the proposed system achieves its intended objectives, several enhancements can further improve its functionality and usability. Future work may include integrating Optical Character Recognition (OCR) to support image-based and scanned PDF documents, enabling the system to process a wider variety of academic resources. Support for multilingual note generation may also be incorporated to assist students studying in different languages. \par

The system can be extended with a complete frontend application that allows students to browse available subjects, chapters, and topics through an intuitive user interface. Additional features such as authentication, role-based access control, document management, and administrative dashboards can further improve system usability and security. Future versions may also support multiple Large Language Models and embedding models through configurable adapters, making the system more flexible and adaptable to evolving AI technologies. \par

Furthermore, retrieval quality may be enhanced through advanced chunking strategies, metadata optimization, and hybrid retrieval techniques. These improvements have the potential to increase the accuracy, coverage, and overall effectiveness of AI-generated academic content. \par